\scanpage{5}

un ens de parties de $S$ ordonné par inclusion.

b) \deleted{L'image} Cet ens de parties est stable par intersection de deux parties minorées

\uline{Dém} b) \inserted{sur $\operatorname{Im}(\Phi\to\mathfrak P(S))$} résulte aussitôt de l'existence des $\operatorname{Inf}$ de deux él. minorés de $\Phi$. Pour a), l'injectivité \inserted{et le fait que $\alpha(x)\leq\alpha(y)\Rightarrow x\leq y$,} résulte du

\deleted{\unsure{Re}} \uline{Corollaire 2} Toute facette est le $\operatorname{Sup}$ des \deleted{facettes} sommets qu'elle majore.

Pour sommet c'est trivial, pour arête c'est 4 a), pour face celà résulte de 4 b) et du

\uline{Corollaire 3} Toute face a au moins trois côtés (i.e. a au moins trois sommets).

En effet, s'il n'en avait que deux, les deux côtés auraient les $=$ ens de sommets sans être égaux, contrairement à 4 a).

\uline{NB} Souvent on \unsure{interprète|interprètera|intersection} les arêtes de l'icosaèdre comme des paires d'arêtes i.e. des él. de $\mathfrak P_2(S)$, les faces comme des parties de $S$ à cinq éléments i.e. des él. de $\mathfrak P_5(S)$.

Dualement

\uline{Corollaire} Tte facette est \deleted{intersection} \inserted{$\operatorname{Inf}$} des faces qui la contiennent. L'application qui, à tte facette $x$, associe l'ens des faces qui la contiennent, définit une injection $S\longrightarrow\mathfrak P(F)$, induisant un

% archive sheet number: 4

